\chapter*{术语索引}
\phantomsection
\addcontentsline{toc}{chapter}{术语索引}

参考编号依次表示小节与目（个别情形为习题）。

同态的右（左）伴随：1, 8。

半线性映射的伴随：7, 3。

二次型的克利福德代数：9, 1。

交错（矩阵）：3, 1。

两条半直线的角：10, 3 及习题 6。

仿射平面上两条射线的角：10, 3。

两条直线的角：10, 3 及习题 2。

仿射平面上两条直线的角：10, 3。

两条有向直线的角：10，习题 2。

平面上两个向量的角：10, 3。

维数 \(\geqslant 2:10,3\) 的空间中两个向量的角。

直角：10, 3 及习题 2。

旋转的角：10, 3。

正相似变换的角：10，习题 6。

平角：10, 3 及习题 6。

二次型的型环：8, 3。

维特环：8, 3。

环的反自同构：1, 2。

克利福德代数的主反自同构：9, 1。

反埃尔米特（型）：3, 1。

反埃尔米特（矩阵）：3, 1。

反对称（矩阵）：3, 1。

双线性映射：1, 1。

右（左）退化的双线性映射：1, 1。

非退化双线性映射：1, 1。

与双线性映射相伴的非退化双线性映射：1, 3。

由双线性映射经标量扩张得到的双线性映射：1, 4。

模到其克利福德代数中的典范映射：9, 1。

从埃尔米特自同态所成之集到埃尔米特型所成之集的典范映射：7, 3。

集合 \(\mathfrak{A}\) 到 \(\mathrm{S}^{+}/\mathrm{H}:10,3\) 上的典范映射。

集合 \(\mathfrak{A}\) 到 \(O^{+}:10,3\) 上的典范映射。

与双线性映射相伴的右（左）线性映射：1, 1。

与半双线性型相伴的右（左）半线性映射：1, 6。

关于一个反自同构的右（左）半双线性映射：1, 2。退化的（右、左）半双线性映射：1, 2。非退化半双线性映射：1, 2。与半双线性映射相伴的非退化半双线性映射：1, 3。由半双线性映射经标量扩张得到的半双线性映射：1, 4。与双线性映射相伴的（线性映射）：1, 1。与半双线性型相伴的（半线性映射）：1, 6。与二次型相伴的（双线性型）：3, 4。正交自同构：6, 2。n 个变量的正交自同构：6, 2。克利福德代数的主自同构：9, 1。辛自同构：5, 3。2m 个变量的辛自同构：5, 3。酉自同构：6, 2。n 个变量的酉自同构：6, 2。仿射线性丛的轴：10，习题 16。正交基：6, 1。标准正交基：6, 1。辛基：5, 1。双线性（映射）：1, 1。双线性（型）：1, 6。双模：1, 1。典范：见典范映射与典范同态。二次曲面的中心：6，习题 25 及 10，习题 12。单位圆：10，习题 2。夏尔（关系式）：10, 3。克利福德（代数）：9, 1。克利福德（群）：9, 5。线性丛：10，习题 16。条件 (C)：6, 1。条件 (C′)：6, 1。条件 (T)：4, 2。迷向锥：6，习题 23。共形（群）：10，习题 14。仿射二次曲线：6，习题 25。射影二次曲线：6，习题 23。共轭（仿射线性簇）：6，习题 25。共轭（射影线性簇）：6，习题 23。角的余弦：10, 3。角的余切：10, 3。维特分解：4, 2。右（左）退化的（双线性映射）：1, 1。右（左）退化的（半线性映射）：1, 2。退化（仿射二次曲线）：6，习题 25。退化（射影二次曲线）：6，习题 23。退化（仿射二次曲面）：6，习题 25。退化（射影二次曲面）：6，习题 23。半直线：10, 3。闭半直线：10, 3。迷向半直线：10, 3。开半直线：10, 3。位移：6, 6。二次型的维数：8, 1。正（相似变换）：6, 5。半双线性型关于一组元素的判别式：2。距离：7, 1。直（角）：10, 3。带基点直线：10，习题 2。

克利福德代数的奇元素：9, 1。迷向元素：4, 1。克利福德代数的偶元素：9, 1。奇异元素：4, 1。关于半双线性映射正交的元素：1, 3。关于二次型正交的元素：3, 4。埃尔米特自同态：7, 3。正规自同态：7, 3。ε-埃尔米特（型）：3, 1。ε-埃尔米特（矩阵）：3, 1。等价（二次型）：3, 4。等价（半双线性型）：1, 6。二次型的定义空间：8, 1。欧几里得空间：6, 6。埃尔米特空间：6, 6。定向空间：10, 3。欧几里得（空间）：6, 6。半双线性型到张量幂的扩张：1, 9。半双线性型到外幂的扩张：1, 9。二次型的外（直和）：3, 4。二次模的外（直和）：3, 4。弱正交（子空间）：3，习题 11。闭（角域）：10, 4。闭（半直线）：10, 3。余弦函数：10, 3。余切函数：10, 3。正弦函数：10, 3。正切函数：10, 3。三角函数：10, 3。反埃尔米特型：3, 1。双线性型：1, 6。与二次型相伴的双线性型：3, 4。埃尔米特双线性型：3, 1。埃尔米特型：3, 1。负（正）埃尔米特型：7, 1。双线性（半双线性）型的逆型：1, 7。度量型：6, 6。中性型：4, 2 及 8, 1。二次型：3, 4。负（正）二次型：7, 1。非退化二次型：3, 4。由二次型经标量扩张得到的二次型：3, 4。半双线性型：1, 6。等价的二次型：3, 4。等价的半双线性型：1, 6。格拉姆–施密特（正交化方法）：6, 1。克利福德群：9, 5。约化克利福德群：9, 5。特殊克利福德群：9, 5。旋转群：9, 5。二次型的型群：8, 2。维特群：8, 2。与 Q 相伴的正交群：6, 2。n 个变量的正交群：6, 2。特殊正交群：6, 2。特殊酉群：6, 2。与 Φ 相伴的辛群：5, 3。2m 个变量的辛群：5, 3。与 Φ 相伴的酉群：6, 2。

\begin{Shaded}
\begin{Highlighting}[]
\NormalTok{n 个变量的酉群 : 6, 2.}
\NormalTok{埃尔米特（自同态） : 7, 3.}
\NormalTok{埃尔米特（空间） : 6, 6.}
\NormalTok{埃尔米特（型） : 3, 1.}
\NormalTok{埃尔米特（矩阵） : 3, 1.}
\NormalTok{\textquotesingle{}典范同态\textquotesingle{}——从 E 的子{-}模的克利福德代数到 E 的克利福德代数中 : 9, 1.\textquotesingle{}}
\NormalTok{度量同态 : 4, 3.}
\NormalTok{两球面的根超平面 : 10, 习题 12.}
\NormalTok{双线性映射的原像 : 1, 1.\textquotesingle{}}
\NormalTok{半双线性映射的原像 : 1, 2.\textquotesingle{}}
\NormalTok{克利福德代数中的奇（元素） : 9, 1.}
\NormalTok{二次型（半双线性型）的指数 : 4, 2.\textquotesingle{}}
\NormalTok{迪克森不变量 : 9, 习题 9.}
\NormalTok{逆（型） : 1, 7.}
\NormalTok{逆（相似变换） : 6, 5.}
\NormalTok{反演 : 10, 习题 13.}
\NormalTok{球面 C 的反演 : 10, 习题 13.}
\NormalTok{线性群中的对合 : 6, 3.}
\NormalTok{迷向（半{-}直线） : 10, 3.}
\NormalTok{迷向（元素） : 4, 1.}
\NormalTok{迷向（子{-}模） : 4, 1.}
\NormalTok{迷向（线性簇） : 6, 6.}
\NormalTok{拉盖尔（公式） : 10, 习题 5.}
\NormalTok{惯性定律 : 7, 2.\textquotesingle{}}
\NormalTok{向量的长度 : 10, 3.\textquotesingle{}}
\NormalTok{交错矩阵 : 3, 1.}
\NormalTok{反埃尔米特矩阵 : 3, 1.}
\NormalTok{反对称矩阵 : 3, 1.}
\NormalTok{满足 (G)、(D)、(GD) 或 (DG) 的映射的矩阵 : 1, 10.\textquotesingle{}}
\NormalTok{双线性映射的矩阵 : 1, 1.\textquotesingle{}}
\NormalTok{半双线性映射的矩阵 : 1, 2.\textquotesingle{}}
\NormalTok{ε{-}埃尔米特矩阵 : 3, 1.}
\NormalTok{埃尔米特矩阵 : 3, 1.}
\NormalTok{正规矩阵 : 7, 习题 17.}
\NormalTok{正交矩阵 : 6, 2.}
\NormalTok{对称矩阵 : 3, 1.}
\NormalTok{辛矩阵 : 5, 3.}
\NormalTok{酉矩阵 : 6, 2.}
\NormalTok{度量（型） : 6, 6.}
\NormalTok{度量（同态） : 4, 3.}
\NormalTok{二次模 : 3, 4.}
\NormalTok{相似变换的乘子 : 6, 5 及 6.\textquotesingle{}}
\NormalTok{负（n{-}向量） : 10, 3.}
\NormalTok{负（埃尔米特型） : 7, 1.}
\NormalTok{负（二次型） : 7, 1.}
\NormalTok{中性（二次型） : 8, 1.}
\NormalTok{中性（半双线性型） : 4, 2.}
\NormalTok{非退化（双线性映射） : 1, 1.}
\NormalTok{非退化（半双线性映射） : 1, 2.}
\NormalTok{非退化（二次型） : 3, 4.}
\NormalTok{正规（自同态） : 7, 3.}
\NormalTok{旋量范数 : 9, 5.}
\NormalTok{二次模的核 : 3, 4.\textquotesingle{}}
\NormalTok{负（正）n{-}向量 : 10, 3.}
\NormalTok{空间的定向 : 10, 3.\textquotesingle{}}
\NormalTok{定向（空间） : 10, 3.}
\NormalTok{正交（自同构） : 6, 2.}
\NormalTok{正交（群） : 6, 2.}
\NormalTok{正交（特殊群） : 6, 2.}
\end{Highlighting}
\end{Shaded}

正交（投影算子）：6, 3。与子模正交的（子模）：1, 3。与线性簇正交的（向量）：6, 6。正交（基）：6, 1。正交（矩阵）：6, 2。到线性簇上的正交（投影）：6, 6。正交（变换）：6, 2。与线性簇正交的（簇）：6, 6。正交（部分）：1, 3 及 3, 4。正交（线性簇）：6, 6。格拉姆–施密特正交化（方法）：6, 1。正交（元素）：1, 3 及 3, 4。标准正交（基）：6, 1。开（角域）：10, 4。开（半直线）：10, 3。克利福德代数中的偶（元素）：9, 1。正交部分：1, 3 及 3, 4。垂直（线性簇）：6，习题 22。交错矩阵的普法夫式：5, 2。平（角）：10, 3。球极投影的视点：10，习题 14。线性簇关于二次曲面的极簇：6，习题 23 及 25。超平面关于二次曲面的极点：6，习题 23 及 25。反演的极点：10，习题 13。正（n-向量）：10, 3。正（埃尔米特型）：7, 1。正（二次型）：7, 1。主（反自同构）：9, 1。主（自同构）：9, 1。二次型的张量积：8, 3。半双线性型的张量积：1, 9。正交投影算子：6, 3。到线性簇上的正交投影：6, 6。球极投影：10，习题 14。伪判别式：9，习题 9。点关于球面的幂：10，习题 12。反演的幂：10，习题 13。勾股（定理）：6, 6。二次（型）：3, 4。二次（模）：3, 4。仿射二次曲面：6，习题 25。射影二次曲面：6，习题 23。双线性（半双线性）型的秩：1, 6。球面的半径：10，习题 12。约化（克利福德群）：9, 5。夏尔关系式：10, 3。旋量表示：9, 4。旋转：9, 5。角为 φ 的旋转：10, 3 及习题 2。旋转（群）：9, 5。闭角域：10, 4。开角域：10, 4。平角域：10，习题 8。凹角域：10，习题 8。凸角域：10，习题 8。半旋量：9, 4。半双线性（映射）：1, 2。半双线性（型）：1, 6。埃尔米特型的符号差：7, 2。相似变换：6, 5 及 6。

正相似变换：6, 5 及 9，习题 9。

逆相似变换：6, 5。

奇异（元素）：4, 1。

奇异（子模）：4, 1。

角的正弦：10, 3。

双线性（半双线性）映射的直和：1, 3。

二次型（二次模）的外直和：3, 4。

迷向子模：4, 1。

与子模正交的子模：1, 3。

奇异子模：4, 1。

全迷向子模：4, 1。

全奇异子模：4, 1。

球面：10，习题 12。

正交球面：10，习题 12。

旋量：9, 4。

半直线的正向序列：10, 4。

关于超平面的对称（变换）：6, 4。

关于向量子空间的对称（变换）：6, 3。

关于仿射线性簇的对称（变换）：6, 6。

对称（矩阵）：3, 1。

辛（自同构）：5, 3。

辛（基）：5, 1。

辛（群）：5, 3。

辛（矩阵）：5, 3。

辛（变换）：5, 3。

角的正切：10, 3。

与二次曲面相切的（线性簇）：6，习题 23 及 25。

勾股定理：6, 6。

维特定理：4, 3。

全迷向（子模）：4, 1。

与子模完全正交的（子模）：1, 3。

与线性簇完全正交的（簇）：6, 6。

全奇异（子模）：4, 1。

正交变换：6, 2。

辛变换：5, 3。

酉变换：6, 2。

三角（函数）：10, 3。

二次型的型：8, 1。

酉（自同构）：6, 2。

酉（群）：6, 2。

酉（特殊群）：6, 2。

酉（矩阵）：6, 2。

酉（变换）：6, 2。

迷向线性簇：6, 6。

与线性簇正交的线性簇：6, 6。

全迷向线性簇：6, 6。

正交的线性簇：6, 6。

与线性簇正交的向量：6, 6。

维特（环）：8, 3。

维特（分解）：4, 2。

维特（群）：8, 2。

维特（定理）：4, 3。

